\section{Síntesis y ampliación}

\begin{enumerate}
    \item RL e IL tienen ventajas distintas; los sistemas reales suelen combinarlos
    \item La sim-to-real transfer es el flujo de trabajo central para aplicar RL a robots reales
    \item La domain randomization y el marco teacher-student son técnicas fundamentales para cerrar la sim-to-real gap
    \item RL ya logró avances decisivos en locomoción con patas y en manipulación diestra
    \item Los robots humanoides representan el reto de la siguiente generación
    \item La calidad de los simuladores y el diseño de recompensas siguen siendo el cuello de botella de las aplicaciones reales
\end{enumerate}

\begin{knowledgebox}{Lecturas adicionales}
\begin{itemize}
    \item Kumar et al., ``Rapid Motor Adaptation (RMA),'' RSS 2021
    \item Tobin et al., ``Domain Randomization for Transferring Deep Neural Networks from Simulation to the Real World,'' IROS 2017
    \item OpenAI, ``Solving Rubik's Cube with a Robot Hand,'' 2019
    \item Rudin et al., ``Learning to Walk in Minutes Using Massively Parallel Deep RL,'' CoRL 2022
    \item Radosavovic et al., ``Real-World Humanoid Locomotion with RL,'' Science Robotics 2024
\end{itemize}
\end{knowledgebox}

\end{document}
